\documentclass[11pt]{article}
\usepackage{amsmath,amssymb}
\usepackage[hidelinks]{hyperref}
\title{Receipt expiry and throughput at a shared custody port}
\date{}
\begin{document}
\maketitle

Two civil custody scopes share a port that can perform one irreversible
handoff per tick. A controller can know the maximum rate of each scope in
isolation and still not know whether the two rates can coexist. In a
matched example they can coexist when one scope may use either of two
admissible phases. Restricting that scope's mechanics, without changing
either isolated rate, makes coexistence impossible. The missing input is
the spacing of the handoffs within one cycle, not the average rate.

\section{Physical interface}
This is a slotted model with a trusted station clock, a saturated offered
workload and a serial custody transfer port. One port use completes a
physical transfer and is the unique mint event for its receipt. Custody
and the authorized signer then change. Old copies are authentic, are
scoped to their own transfer stream, can arrive through \(H\) ticks
after mint inclusive, and cannot be minted again or age-reset. The
earliest safe reuse of the same signer and tag is \(W=H+1\) ticks
after a mint. Assume admissible saturated carrier work is available
whenever the stated mechanical minimum is reached; arbitrary delay of
clearance, transfer or observation is not covered. Receipt freshness
depends only on scope, signer and tag, and a surviving stale copy is
allowed without affecting the prior local history.

This specializes the recurrence in \#388, in which expiry is measured
by local advances, physical transfer and separate observations are
different events, and a shared port may be only a receiving port. Here
clock, unique mint and one occupied slot per handoff in both streams
are explicit extra plant assumptions. There is no promise of a
deadline for each offered job.

\section{Isolated openings and sharp tag reuse}
A scope has two alternating authenticated signers. Consecutive transfers
require at least \(d_0,d_1>0\) slots, indexed by the \emph{next} signer.
Write \(D=d_0+d_1\). For \(M\) tag values per signer, set \(K=2M\).
The necessary and sufficient reuse constraint for any schedule
\(t_0<t_1<\cdots\) is
\[
       t_g-t_{g-K}\geq W \qquad (g\geq K).
\]
For necessity, if this interval is shorter than \(W\), the \(M+1\)
same-signer mints in it form a clique of possible authentic stale
conflicts; a pair must reuse a tag. For sufficiency, assign epoch \(g\)
tag \(\lfloor g/2\rfloor\bmod M\); only the last mint of each pair
matters because mint and expiry order coincide. The lower bound is
for the stated receipt-only decoder, not for protocols that can obtain
other authoritative evidence of the current clearance.

On one stream, the componentwise earliest safe sequence starts at
zero and chooses the maximum of the minimum physical next time and
\(t_{g-K}+W\) when defined. Let \(P=\max(MD,W)\). Writing \(g=qK+r\)
with \(0\leq r<K\), the result is
\[
       t_g=qP+\lfloor r/2\rfloor D+(r\bmod 2)d_1.
\]
Inside a block the physical successor time equals this expression;
at the wrap the recurrence picks the larger of the physical block
minimum and the expiry wait. This also shows that the isolated rate
is \(K/P\). The smallest unpaced alphabet is \(\lfloor H/D\rfloor+1\).
A larger elapsed age can lower the rate even if the same clockless
advance-count certificate appeared adequate.

\section{Exact coexistence of isolated saturation rates}
For each scope form all \(K\)-point calendars in one cycle of length
\(P\), with first signer 0 at position zero, strictly increasing
positions, and every cyclic gap at least the mechanical minimum for
its successor signer. For two calendars \(S_A,S_B\), a pair of shifted
periodic streams is collision-free iff the prohibited offsets
\[
         \{s_A-s_B\bmod \gcd(P_A,P_B):s_A\in S_A,\ s_B\in S_B\}
\]
do not fill all residue classes modulo that greatest common divisor.
This is the elementary intersection test for two arithmetic
progressions. Existence of such calendars is necessary and sufficient
for simultaneously attaining the two isolated rates, \emph{even if the
shared-port controller may choose arbitrary nonperiodic schedules}.

For necessity, the port model is a finite automaton: retain next
phase and the ages, capped above their largest consumer's threshold,
of the last \(K\) port uses of each scope. Invalid actual
effects are excluded on every transition. Each cycle has both rates
at most their isolated maxima. Finite paths decompose as cycles and a
bounded remainder. If no cycle attained both individual maxima, the
sum of normalized rates would be uniformly less than two on the finite
set of simple cycles, so no nonperiodic history could attain both.
In a cycle attaining both, consider a scope's successive \(K\)-event
lags. Each lag is at least \(P\) by the mechanical and reuse
constraints. Their average equals \(P\) at the isolated rate, hence
all are exactly \(P\). This gives the asserted calendar, including
when only the mechanical bound is active. Sufficiency follows by
running the schedules with an admissible offset at the same port,
starting each scope's first cycle only when its signer 0 is due.
Initial unused receipt slots cannot introduce predecessor obstacles.

With one tag, each scope has two residues separated by
\(d_1\leq\delta\leq P-d_0\). If the common divisor is two, both scopes
must choose even \(\delta\), otherwise one uses both residues.
Take \(B=(d_0,d_1,W,M)=(1,1,6,1)\) and compare
\(A=(1,1,4,1)\) with \(A'=(1,3,4,1)\). Both pairs have isolated rates
\(1/2\) and \(1/3\). For the first, use even gaps of two and opposite
phases; both scopes can transfer at their isolated rates.
For the second, the changed local preparation law forces an odd gap
of three and prevents that. Reopening earlier by disregarding the
mechanical gap or stale-copy support is not a repair.

For the second pair the complete limiting rate region for
nonnegative rates \(x,y\) is
\[
     x\leq \tfrac12,\qquad y\leq\tfrac13,\qquad 3x+2y\leq2 .
\]
The certificate gives an integer potential on all 238 reachable
states. On each of the 498 transitions its increment dominates
\(3\,1_A+2\,1_B-2\). Summing bounds every finite prefix up to a
constant, therefore gives the third inequality on limiting rates.
The first two are the isolated bounds. Two periodically repeated
schedules achieve \((1/2,1/4)\) and \((4/9,1/3)\), and the other
vertices are obtained by running either scope alone or idling.
Convex combinations are approached by concatenating increasingly
long periods with bounded idle guards; each scope can finish a block
on its original signer, and after a guard the old receipt ages no
longer constrain the next block. A sequence of blocks growing slowly
enough relative to elapsed time gives the claimed limiting rates,
not merely a subsequence of endpoint averages. The bound has an
independent integer checker that also replays the schedules using
uncapped physical times. No solver tolerance is used.

\section{When the mint is late}
The assessment changes if a receipt can instead be minted after the
prior nominal opening. The relevant date for first reuse is the
latest physically possible mint in the received prior-history fibre
plus \(W\), not the old opening plus \(W\). For example, with \(H=8\),
one tag and earlier completion in three ticks, evidence that certified
mint no later than tick two permits reuse at tick eleven; a bound
of seven from an unresolved earlier phase only supports reuse at
tick sixteen. A genuinely late prior mint at seven makes sixteen
necessary if that copy can survive independently. These are event
bounds, not a statement that an actor receives that evidence in time.
The exact calendar result above assumes mint at the shared port use;
it does not feed a late or schedule-dependent mint back into its theorem.

\section{Comparators and status}
\href{https://www.rfc-editor.org/rfc/rfc9293.html#section-3.4.2}{RFC 9293
\(\S3.4.2\)} supplies a practical comparison for enforcing a
maximum duplicate lifetime before sequence reuse. The strictly
periodic collision part is standard scheduling; compare
\href{https://drops.dagstuhl.de/storage/16dagstuhl-seminar-proceedings/dsp-vol10071/DagSemProc.10071.13/DagSemProc.10071.13.pdf}{Eisenbrand
et al.}. The source argument explains when that comparator is exact
for a controller that initially permits arbitrary port histories.
Finite model checks compare the calendar criterion with direct
weighted-cycle search on the same information and plant and tie.
The point is the sufficiency or insufficiency of a source interface
that exports only throughput, not a new scheduling algorithm.
This is a restricted recurrence result, not unbounded decentralized
controller synthesis nor acceptance of any integration obligation.
\end{document}
